\documentclass[11pt]{article}
\usepackage[margin=1in]{geometry}
\usepackage[T1]{fontenc}
\usepackage{graphicx}
\usepackage{hyperref}
\usepackage{enumitem}
\setlength{\parskip}{0.5em}
\setlength{\parindent}{0pt}

\title{large language model agents and tool use}
\author{arXiv GraphRAG --- phase-6 deep research}
\date{\today}

\begin{document}
\maketitle

\begin{abstract}
The review is exceptionally well-structured and covers the text/API frontier deeply, but needs more breadth in multimodal/visual agency and operational orchestration.
\end{abstract}

\section{From Text Prediction to Autonomous Action: A Review of Tool-Augmented Large Language Model Agents}

The rapid evolution of large language models (LLMs) has shifted the computational paradigm from simple probabilistic text prediction toward complex cognitive processing. While the \textit{Foundations of Language Modeling and Reasoning} established the capacity for sophisticated linguistic patterns, inherent limitations such as factual hallucinations and a lack of real-world grounding necessitate a transition toward agency. The ability of a model to act autonomously—to reason, plan, and interact with external environments—represents a critical leap in artificial intelligence. This evolution is vital because it transforms LLMs from passive knowledge repositories into active problem-solving entities capable of executing complex, multi-step tasks by bridging the gap between internal reasoning and external execution.

This review traces the trajectory of agentic development, beginning with \textit{Early Paradigms} centered on reasoning, planning, and retrieval-augmented generation. It then examines the pivotal shift toward \textit{The Emergence of Tool-Augmented Language Models}, where models began leveraging external software and APIs to extend their functional capabilities beyond the confines of their training data. The discussion expands into the structural dimensions of AI through \textit{Multi-Agent Systems and Collaborative Frameworks}, alongside an analysis of \textit{Comprehensive Surveys and Agentic Methodologies}. To address the critical challenges of reliability and performance, the review evaluates \textit{Advanced Benchmarking and Evaluation Frameworks} before concluding at \textit{The Frontier: Agentic Reasoning and Orchestration}, where the focus shifts toward the seamless integration of autonomous reasoning with complex, high-level task orchestration.

\subsection{Foundations of Language Modeling and Reasoning}

The evolution of language modeling has transitioned from fundamental architectural innovations to complex reasoning and alignment strategies. The foundation of this progression is the Transformer architecture (\texttt{arXiv:1706.03762}), which replaced traditional recurrence and convolution with attention mechanisms. This architecture proved highly effective for machine translation, achieving BLEU scores of 28.4 on English-to-German and 41.8 on English-to-French tasks (\texttt{arXiv:1706.03762}). Following this, research diverged into two primary directions: bidirectional understanding and autoregressive generation. BERT (\texttt{arXiv:1810.04805}) utilized deep bidirectional representations to excel at language understanding, yielding a GLUE score of 80.5\% and high SQuAD v1.1 Test F1 scores of 93.2. In contrast, GPT-3 (\texttt{arXiv:2005.14165}) leveraged an autoregressive architecture with 175 billion parameters to enable few-shot learning. Unlike BERT’s focus on context-heavy understanding, GPT-3 allows for text-based interaction via few-shot demonstrations without requiring gradient updates or fine-tuning (\texttt{arXiv:2005.14165}). However, GPT-3 faces challenges regarding methodological issues stemming from training on large web corpora and potential societal impacts (\texttt{arXiv:2005.14165}).

As models grew in scale, the focus shifted from simple text completion to eliciting complex reasoning and ensuring human alignment. Chain-of-thought prompting (\texttt{arXiv:2201.11903}) introduced a mechanism to improve performance on arithmetic, commonsense, and symbolic reasoning tasks—such as GSM8K math word problems—by generating intermediate reasoning steps. This differs from the standard few-shot prompting used in GPT-3 (\texttt{arXiv:2005.14165}) by explicitly structuring the model's cognitive process. To bridge the gap between raw model capabilities and user intent, the two-stage fine-tuning process of supervised fine-tuning (SFT) followed by reinforcement learning from human feedback (RLHF) was developed (\texttt{arXiv:2203.02155}). While this approach enhances truthfulness and reduces toxicity, it is not without flaws; models still make simple mistakes and may experience minimal performance regressions on certain public NLP datasets (\texttt{arXiv:2203.02155}).

Recent developments suggest that these next-word prediction models are approaching the threshold of artificial general intelligence. GPT-4 demonstrates significant capabilities across diverse domains, including mathematics, coding, medicine, and law (\texttt{arXiv:2303.12712}). Despite these "sparks" of AGI, the system remains incomplete (\texttt{arXiv:2303.12712}). A critical debate persists as to whether the current trajectory of next-word prediction is sufficient, or if a new paradigm beyond this framework is required to achieve true intelligence (\texttt{arXiv:2303.12712}).

\subsection{Early Paradigms: Reasoning, Planning, and Retrieval}

Early research into LLM agents focused on bridging the gap between static linguistic generation and dynamic, goal-oriented execution through structured reasoning and planning. A foundational approach is the ReAct paradigm, which interleaves reasoning traces with task-specific actions to induce, track, and update action plans [\texttt{arXiv:2210.03629}]. By utilizing external tools such as Wikipedia APIs and knowledge bases, ReAct enables agents to perform question answering and fact verification through a synergistic loop of thought and action [\texttt{arXiv:2210.03629}].

While ReAct emphasizes the interleaving of thought and action, other works focus on the decomposition of high-level tasks into executable sequences. For instance, zero-shot planning frameworks attempt to decompose natural language instructions into mid-level plans, which are then translated into admissible actions via semantic translation [\texttt{arXiv:2201.07207}]. However, this approach faces a significant trade-off between correctness and executability, as naive LLM plans often fail to map precisely to available actions [\texttt{arXiv:2201.07207}]. This challenge of grounding is further addressed in the "Inner Monologue" framework, which employs a closed-loop reasoning mechanism [\texttt{arXiv:2207.05608}]. Unlike purely generative planners, Inner Monologue uses natural language feedback from the environment—such as scene descriptions or success detection—to inform continuous decision-making in embodied tasks like kitchen manipulation [\texttt{arXiv:2207.05608}].

The literature also reveals a tension between linear reasoning and structured exploration. Standard LLM inference is often limited by left-to-right, token-level decision-making, which is insufficient for tasks requiring strategic lookahead [\texttt{arXiv:2305.10601}]. To overcome this, the Tree of Thoughts (ToT) framework generalizes reasoning by exploring multiple coherent "thoughts" through self-evaluation and backtracking [\texttt{arXiv:2305.10601}]. This structured exploration yields significant performance gains; for example, in the Game of 24, ToT achieved a 74\% success rate compared to only 4\% for GPT-4 using Chain-of-Thought [\texttt{arXiv:2305.10601}]. Similarly, the Reasoning via Planning (RAP) framework repurposes the LLM as both a reasoning agent and a world model, utilizing Monte Carlo Tree Search (MCTS) to balance exploration and exploitation [\texttt{arXiv:2305.14992}]. RAP addresses the inherent limitation that LLMs lack an internal world model to simulate long-term outcomes, achieving a 33\% relative improvement over GPT-4's CoT in plan generation [\texttt{arXiv:2305.14992}].

Finally, agents have moved toward self-improvement through iterative feedback. The Reflexion framework introduces verbal reinforcement learning, where agents maintain an episodic memory buffer of linguistic feedback to improve subsequent trials [\texttt{arXiv:2303.11366}]. While this allows for high performance in coding tasks—achieving 91\% pass@1 on HumanEval [\texttt{arXiv:2303.11366}]—it highlights a persistent challenge: the difficulty of learning efficiently from trial-and-error without extensive training or expensive fine-tuning [\texttt{arXiv:2303.11366}].

\subsection{The Emergence of Tool-Augmented Language Models}

The evolution of tool-augmented language models has transitioned from specialized reasoning offloading to the management of massive, multi-modal API ecosystems. Early efforts focused on mitigating the inherent logical and arithmetic errors found in standard chain-of-thought prompting [\texttt{arXiv:2211.10435}]. For instance, Program-Aided Language models (PAL) [\texttt{arXiv:2211.10435}] decompose natural language problems into runnable Python programs, offloading execution to a symbolic interpreter to ensure mathematical and algorithmic accuracy. In contrast, Toolformer [\texttt{arXiv:2302.04761}] adopts a self-supervised learning approach, teaching the model to decide when and how to call various APIs—such as calculators, search engines, or translation systems—by incorporating tool results directly into its token prediction process.

As the scope of tool use expanded, research shifted toward complex task planning and multi-modal integration. HuggingGPT [\texttt{arXiv:2303.17580}] introduces an LLM-powered controller that manages a diverse suite of AI models from the Hugging Face community, spanning language, vision, and speech modalities. While HuggingGPT focuses on high-level task planning and model selection via a generic language interface, other works address the technical precision required for API interaction. Gorilla [\texttt{arXiv:2305.15334}] specifically targets the tendency of standard LLMs to hallucinate incorrect API usage or arguments. By combining a fine-tuned LLaMA-based model with a document retriever, Gorilla improves the accuracy of API calls for platforms like HuggingFace, TorchHub, and TensorHub.

The scale of available tools has since grown exponentially, moving beyond specialized libraries to real-world web services. ToolLLM [\texttt{arXiv:2307.16789}] pushes this boundary by mastering over 16,464 real-world RESTful APIs across 49 categories from the RapidAPI Hub. Unlike the retrieval-augmented approach of Gorilla [\texttt{arXiv:2305.15334}], ToolLLM utilizes a neural API retriever and a depth-first search-based decision tree to facilitate complex, multi-tool reasoning traces and solution paths.

Despite these advancements, the field faces structural challenges. A general framework for tool learning [\texttt{arXiv:2304.08354}] emphasizes the necessity of task decomposition and dynamic plan adjustment to conquer sub-tasks. However, this work also highlights significant gaps in the current literature, noting a lack of systematic tool-learning evaluation in prior works and a deficiency in the comprehensive understanding of key challenges and future opportunities within the domain [\texttt{arXiv:2304.08354}].

\subsection{Multi-Agent Systems and Collaborative Frameworks}

The evolution of large language model (LLM) applications has shifted from single-agent interactions toward sophisticated multi-agent collaborative frameworks. Early approaches, such as the CAMEL framework, utilize an inception prompting mechanism to facilitate autonomous cooperation through role-playing agents (\texttt{arXiv:2303.17760}). While CAMEL establishes a foundation for instruction-following cooperation, it remains limited by its reliance on human input to guide the conversational trajectory (\texttt{arXiv:2303.17760}).

Building upon these role-playing foundations, subsequent research has sought to introduce more rigorous structural constraints to improve task performance. MetaGPT moves away from open-ended role-play toward an "assembly line" paradigm, where Standardized Operating Procedures (SOPs) are encoded into prompt sequences to mimic human professional workflows, particularly in software engineering (\texttt{arXiv:2308.00352}). However, the efficiency of this structured approach is tempered by the risk of logic inconsistencies caused by cascading hallucinations when LLMs are naively chained together (\texttt{arXiv:2308.00352}). In contrast to the rigid SOPs of MetaGPT, AutoGen offers a more flexible conversational architecture. It enables customizable, conversable agents that can interact via natural language or computer code, and importantly, allows for the integration of human inputs and external tools to support tasks ranging from mathematics to online decision-making (\texttt{arXiv:2308.08155}).

While frameworks like AutoGen and MetaGPT focus on workflow and customization, other research explores the cognitive benefits of agent interaction to improve reasoning. The "society of minds" approach in multi-agent debate involves multiple rounds of proposing and debating individual responses and reasoning processes to reach a consensus, specifically aimed at improving factual validity and mathematical reasoning (\texttt{arXiv:2305.14325}). This iterative refinement is echoed in the peer-review collaboration strategy, which emulates academic processes by having agents construct solutions, provide reviews, assign confidence levels, and ultimately revise their work based on received feedback (\texttt{arXiv:2311.08152}).

Finally, the focus of multi-agent research has expanded to include the study of emergent social properties. AgentVerse introduces a framework inspired by human group dynamics, allowing for the dynamic composition of agent groups to accomplish general tasks (\texttt{arXiv:2308.10848}). While this approach facilitates the exploration of emergent social behaviors, it introduces the unique challenge of mitigating negative social behaviors while leveraging positive ones to enhance collective potential (\texttt{arXiv:2308.10848}). Together, these frameworks represent a spectrum of collaboration, ranging from highly structured professional workflows to fluid, debate-driven, or socially dynamic interactions.

\subsection{Comprehensive Surveys and Agentic Methodologies}

The evolution of Large Language Model (LLM) agents is characterized by a transition from general architectural frameworks to specialized methodologies for tool integration and complex task planning. Early foundational surveys established a tripartite agent architecture consisting of a "brain," perception, and action components, categorizing applications into single-agent, multi-agent, and human-agent cooperation scenarios [\texttt{arXiv:2309.07864}]. While this framework provides a broad structural baseline, it faces the critical limitation of lacking a powerful, general model capable of adapting to diverse, unpredictable environments [\texttt{arXiv:2309.07864}].

Recent literature has moved toward more granular taxonomies of agentic behavior. For instance, TPTU introduces a structured framework that distinguishes between one-step and sequential agents to address complex tasks requiring both planning and tool usage [\texttt{arXiv:2308.03427}]. This focus on the intersection of planning and execution is echoed in more recent surveys, which define agentic interaction through mechanisms such as prompting, reasoning, tool utilization, and in-context learning [\texttt{arXiv:2404.04442}]. Specifically, the paradigm for tool-augmented agents has matured into a functional pipeline: understanding user intent, selecting appropriate tools, and performing dynamic plan adjustment [\texttt{arXiv:2409.18807}].

A significant point of divergence in the literature concerns the depth of agentic autonomy and the complexity of their internal cognitive processes. Some frameworks emphasize the integration of cognitive mechanisms—such as reasoning, planning, and memory—within both single and multi-agent systems [\texttt{arXiv:2508.17281}]. Others adopt a methodology-centered taxonomy that links architectural foundations directly to evolutionary pathways and collaboration mechanisms [\texttt{arXiv:2503.21460}]. This methodological shift reflects a growing recognition that agentic success depends not just on the LLM's generative capacity, which may be insufficient for complex tasks [\texttt{arXiv:2308.03427}], but on the robustness of the underlying reasoning and planning loops.

Evaluation methodologies have similarly expanded. While early work focused on specific task planning and tool usage (TPTU) abilities [\texttt{arXiv:2308.03427}], contemporary reviews analyze a wide array of benchmarks, including AgentBench, WebArena, and ToolLLM [\texttt{arXiv:2404.04442}], or aggregate findings across dozens of publicly available datasets [\texttt{arXiv:2508.17281}]. Despite these advancements, several persistent challenges remain central to the discourse. These include the accuracy of tool selection and invocation timing [\texttt{arXiv:2409.18807}], the mitigation of hallucinations [\texttt{arXiv:2404.04442}], and the difficulty of achieving verifiable reasoning and self-improvement in autonomous agents [\texttt{arXiv:2508.17281}]. Consequently, the field is currently navigating the tension between expanding the functional scope of agents—such as moving from tool users to tool creators [\texttt{arXiv:2409.18807}]—and the need for better human value alignment and personalization [\texttt{arXiv:2404.04442}, \texttt{arXiv:2508.17281}].

\subsection{Advanced Benchmarking and Evaluation Frameworks}

As large language model (LLM) agents transition from simple text generation to complex tool use, the development of specialized evaluation frameworks has become critical. Current research demonstrates a shift from general capability testing toward highly specialized, domain-specific, and interaction-centric benchmarks.

Early foundational work like AgentBench provides a multi-dimensional approach by evaluating agents across eight distinct interactive environments, utilizing both API-based and open-source (OSS) tool categories to assess reasoning and decision-making (\texttt{arXiv:2308.03688}). While AgentBench establishes a baseline for agentic performance, it reveals significant structural weaknesses, specifically regarding long-term reasoning, instruction following, and a notable performance gap between top commercial models and OSS competitors under 70B parameters (\texttt{arXiv:2308.03688}). This gap is echoed in the GAIA benchmark, which focuses on real-world questions requiring web browsing and multi-modality handling (\texttt{arXiv:2311.12983}). GAIA highlights a stark disparity in general intelligence; while humans achieve 92\% success on conceptually simple tasks, GPT-4 with plugins scores only 15\% (\texttt{arXiv:2311.12983}).

Recent literature has moved toward evaluating agents within dynamic, high-stakes environments. \$\textbackslash{}tau\$-bench shifts the focus to tool-agent-user interaction, simulating real-world retail domains through dynamic conversations between an agent and a simulated user (\texttt{arXiv:2406.12045}). However, \$\textbackslash{}tau\$-bench exposes the fragility of current models, noting that even GPT-4o struggles with domain-specific rules, yielding low success rates (<50\%) and high inconsistency in retail tasks (\texttt{arXiv:2406.12045}). Similarly, MLE-bench evaluates agents on the highly technical task of machine learning engineering, such as training models and preparing datasets, using Kaggle competitions as a baseline (\texttt{arXiv:2410.07095}). While MLE-bench provides a rigorous technical standard, it faces unique challenges regarding resource scaling and the potential for pre-training contamination (\texttt{arXiv:2410.07095}).

The frontier of evaluation is increasingly defined by diagnostic depth and automated generation. The Agent Planning Benchmark (APB) provides a granular diagnostic framework across 22 domains, testing not just success, but plan correctness and robustness against "broken" or "extraneous" tools (\texttt{arXiv:2606.04874}). This addresses the long-horizon planning and tool-noise weaknesses identified in earlier works like AgentBench (\texttt{arXiv:2308.03688}; \texttt{arXiv:2606.04874}). Furthermore, ActuBench introduces a sophisticated multi-agent pipeline—utilizing roles such as a "distractor constructor" and an "independent verifier"—to automate the generation of complex actuarial reasoning tasks (\texttt{arXiv:2604.20273}). While ActuBench offers a scalable way to create difficult assessments, it notes a critical methodological tension: multiple-choice scaffolds may artificially inflate performance ceilings, necessitating the use of LLM-as-Judge evaluation to accurately differentiate performance at the frontier (\texttt{arXiv:2604.20273}).

\subsection{The Frontier: Agentic Reasoning and Orchestration}

The evolution of large language model (LLM) agents is characterized by a shift from simple prompting to complex, multi-layered orchestration involving reasoning, tool integration, and self-evolution. Recent research identifies three primary paradigms for agentic tool use: plug-and-play prompting, supervised tool learning, and reward-driven tool policy learning (\texttt{arXiv:2604.00835}). While these paradigms offer varying degrees of integration, they face persistent challenges regarding task fragmentation and diverse failure modes across different tool types (\texttt{arXiv:2604.00835}).

A critical tension exists between free-form reasoning and structured execution. For instance, while reinforcement learning (RL) can incentivize emergent reasoning patterns—such as self-reflection and dynamic strategy adaptation—in mathematics and STEM fields (\texttt{arXiv:2501.12948}), free-form Chain-of-Thought (CoT) often proves counterproductive for highly structured tasks (\texttt{arXiv:2509.18076}). To mitigate errors in parameterization and tool selection, researchers have proposed curriculum-inspired frameworks using guided-structured templates (\texttt{arXiv:2509.18076}). These templates provide step-by-step guidance for function calling, achieving relative tool-use error reductions of 3-12\% compared to strong baselines (\texttt{arXiv:2509.18076}).

The architectural complexity of agents is further expanding through hierarchical and multi-agent frameworks. One comprehensive three-layered model proposes a progression from foundational reasoning (planning and search) to self-evolving reasoning (feedback and memory) and, finally, to collective multi-agent reasoning involving coordination and knowledge sharing (\texttt{arXiv:2601.12538}). This trajectory mirrors the broader evolution of tool use, which has moved from single-tool calls toward multi-tool orchestration over long trajectories, utilizing intermediate state and execution feedback to navigate digital environments, software systems, and mobile interfaces (\texttt{arXiv:2603.22862}).

However, as agents become more autonomous, significant bottlenecks emerge. While autonomous systems can now orchestrate complete benchmark construction pipelines—managing everything from subtask design to quality control (\texttt{arXiv:2606.06462})—these processes remain labor-intensive and face sustainability concerns (\texttt{arXiv:2606.06462}). Furthermore, the transition to real-world deployment is hindered by several unresolved challenges: the need for better world modeling, long-horizon interaction capabilities, personalization, and the governance of multi-agent systems (\texttt{arXiv:2601.12538}). Additionally, as orchestration scales, issues regarding safety, cost, verifiability, and reliability become increasingly paramount (\texttt{arXiv:2603.22862}). Together, these developments suggest that while agentic reasoning is maturing through RL and structured templates, the frontier lies in overcoming the scalability and reliability gaps inherent in complex, multi-tool environments.

\subsection{Tensions and Open Problems}

A primary debate in the literature concerns the optimal mechanism for eliciting complex reasoning: whether it is more effective to rely on internal linguistic chain-of-thought processes or to offload computation to external symbolic engines. While early work suggested that self-consistency in linguistic reasoning could significantly improve accuracy (\texttt{arXiv:2203.11171}), subsequent research argues that zero-shot reasoning techniques can achieve comparable results without explicit prompting (\texttt{arXiv:2205.11916}). This tension extends to the use of code, where a divide exists between models that use natural language to reason and those that utilize programmatic execution. Specifically, while Program of Thoughts (PoT) attempts to disentangle computation from language (\texttt{arXiv:2211.12588}), the Program-aided Language Model (PAL) approach argues that direct program generation provides a more robust framework for solving mathematical and logical tasks (\texttt{arXiv:2211.10435}).

Another significant tension involves the architectural approach to agentic autonomy and planning. Some frameworks advocate for interactive, multi-step planning cycles where agents describe, explain, and select actions based on environmental feedback (\texttt{arXiv:2302.01560}). However, this iterative approach has been critiqued for its potential inefficiency compared to more structured reinforcement learning methods like Reflexion (\texttt{arXiv:2303.11366}), or for failing to achieve the high-fidelity social simulation seen in generative agent architectures (\texttt{arXiv:2304.03442}). Furthermore, the debate over "optimal" planning persists; while interactive planning offers flexibility, the LLM+P framework suggests that integrating LLMs with formal optimal planning algorithms provides a more reliable path to task completion than purely linguistic planning (\texttt{arXiv:2304.11477}).

Finally, there is a growing disagreement regarding the scale and nature of "emergent" capabilities and tool proficiency. While foundational studies suggested that scaling parameters leads to the sudden emergence of new abilities (\texttt{arXiv:2206.07682}), more recent technical reports argue that these abilities are often more predictable continuations of few-shot learning capabilities (\texttt{arXiv:2303.08774}). This skepticism extends to tool use; whereas early benchmarks focused on specialized tool-augmented environments (\texttt{arXiv:2304.08244}, \texttt{arXiv:2306.05301}), newer research argues that these are too narrow to represent true mastery. Instead, the field is moving toward massive-scale tool orchestration (\texttt{arXiv:2307.16789}) and self-taught tool use (\texttt{arXiv:2302.04761}), suggesting that the true frontier lies not in following instructions, but in the autonomous acquisition and orchestration of complex API ecosystems.

\subsection{Outlook}

The field is rapidly transitioning from single-turn tool invocation toward complex, multi-agent orchestration and autonomous "agentic reasoning" (\texttt{arXiv:2601.12538}, \texttt{arXiv:2603.22862}). As the complexity of tasks scales, the focus is shifting from simple prompt engineering to the development of scalable multi-agent environments (\texttt{arXiv:2601.09750}) and rigorous, diagnostic benchmarking frameworks (\texttt{arXiv:2604.06462}, \texttt{arXiv:2606.04874}). The next frontier will likely be defined by the convergence of embodied action and high-level strategic planning, moving beyond mere text generation toward agents capable of navigating real-world, multi-step pipelines with verifiable accuracy and generalizable intelligence.

\section*{References}
\begin{enumerate}\small
  \item Attention Is All You Need (2017). \texttt{arXiv:1706.03762}
  \item Proximal Policy Optimization Algorithms (2017). \texttt{arXiv:1707.06347}
  \item BERT: Pre-training of Deep Bidirectional Transformers for Language
  Understanding (2018). \texttt{arXiv:1810.04805}
  \item HotpotQA: A Dataset for Diverse, Explainable Multi-hop Question
  Answering (2018). \texttt{arXiv:1809.09600}
  \item Language Models are Few-Shot Learners (2020). \texttt{arXiv:2005.14165}
  \item Retrieval-Augmented Generation for Knowledge-Intensive NLP Tasks (2020). \texttt{arXiv:2005.11401}
  \item Measuring Massive Multitask Language Understanding (2020). \texttt{arXiv:2009.03300}
  \item ALFWorld: Aligning Text and Embodied Environments for Interactive
  Learning (2020). \texttt{arXiv:2010.03768}
  \item Evaluating Large Language Models Trained on Code (2021). \texttt{arXiv:2107.03374}
  \item Training Verifiers to Solve Math Word Problems (2021). \texttt{arXiv:2110.14168}
  \item Finetuned Language Models Are Zero-Shot Learners (2021). \texttt{arXiv:2109.01652}
  \item WebGPT: Browser-assisted question-answering with human feedback (2021). \texttt{arXiv:2112.09332}
  \item Measuring Mathematical Problem Solving With the MATH Dataset (2021). \texttt{arXiv:2103.03874}
  \item LoRA: Low-Rank Adaptation of Large Language Models (2021). \texttt{arXiv:2106.09685}
  \item Chain-of-Thought Prompting Elicits Reasoning in Large Language Models (2022). \texttt{arXiv:2201.11903}
  \item Training language models to follow instructions with human feedback (2022). \texttt{arXiv:2203.02155}
  \item PaLM: Scaling Language Modeling with Pathways (2022). \texttt{arXiv:2204.02311}
  \item ReAct: Synergizing Reasoning and Acting in Language Models (2022). \texttt{arXiv:2210.03629}
  \item Do As I Can, Not As I Say: Grounding Language in Robotic Affordances (2022). \texttt{arXiv:2204.01691}
  \item Self-Consistency Improves Chain of Thought Reasoning in Language Models (2022). \texttt{arXiv:2203.11171}
  \item Large Language Models are Zero-Shot Reasoners (2022). \texttt{arXiv:2205.11916}
  \item Language Models as Zero-Shot Planners: Extracting Actionable Knowledge
  for Embodied Agents (2022). \texttt{arXiv:2201.07207}
  \item Inner Monologue: Embodied Reasoning through Planning with Language
  Models (2022). \texttt{arXiv:2207.05608}
  \item Emergent Abilities of Large Language Models (2022). \texttt{arXiv:2206.07682}
  \item WebShop: Towards Scalable Real-World Web Interaction with Grounded
  Language Agents (2022). \texttt{arXiv:2207.01206}
  \item Beyond the Imitation Game: Quantifying and extrapolating the
  capabilities of language models (2022). \texttt{arXiv:2206.04615}
  \item PAL: Program-aided Language Models (2022). \texttt{arXiv:2211.10435}
  \item Self-Instruct: Aligning Language Models with Self-Generated Instructions (2022). \texttt{arXiv:2212.10560}
  \item A Generalist Agent (2022). \texttt{arXiv:2205.06175}
  \item ProgPrompt: Generating Situated Robot Task Plans using Large Language
  Models (2022). \texttt{arXiv:2209.11302}
  \item Program of Thoughts Prompting: Disentangling Computation from Reasoning
  for Numerical Reasoning Tasks (2022). \texttt{arXiv:2211.12588}
  \item TALM: Tool Augmented Language Models (2022). \texttt{arXiv:2205.12255}
  \item Code as Policies: Language Model Programs for Embodied Control (2022). \texttt{arXiv:2209.07753}
  \item TPTU: Large Language Model-based AI Agents for Task Planning and Tool Usage (2023). \texttt{arXiv:2308.03427}
  \item Large Language Models as Tool Makers (2023). \texttt{arXiv:2305.17126}
  \item ModelScope-Agent: Building Your Customizable Agent System with
  Open-source Large Language Models (2023). \texttt{arXiv:2309.00986}
  \item An In-depth Survey of Large Language Model-based Artificial Intelligence
  Agents (2023). \texttt{arXiv:2309.14365}
  \item Towards Reasoning in Large Language Models via Multi-Agent Peer Review
  Collaboration (2023). \texttt{arXiv:2311.08152}
  \item Agent Instructs Large Language Models to be General Zero-Shot Reasoners (2023). \texttt{arXiv:2310.03710}
  \item Controlling Large Language Model-based Agents for Large-Scale
  Decision-Making: An Actor-Critic Approach (2023). \texttt{arXiv:2311.13884}
  \item A Survey on Large Language Model based Autonomous Agents (2023). \texttt{arXiv:2308.11432}
  \item GPT-4 Technical Report (2023). \texttt{arXiv:2303.08774}
  \item Toolformer: Language Models Can Teach Themselves to Use Tools (2023). \texttt{arXiv:2302.04761}
  \item LLaMA: Open and Efficient Foundation Language Models (2023). \texttt{arXiv:2302.13971}
  \item Generative Agents: Interactive Simulacra of Human Behavior (2023). \texttt{arXiv:2304.03442}
  \item Reflexion: Language Agents with Verbal Reinforcement Learning (2023). \texttt{arXiv:2303.11366}
  \item Tree of Thoughts: Deliberate Problem Solving with Large Language Models (2023). \texttt{arXiv:2305.10601}
  \item Sparks of Artificial General Intelligence: Early experiments with GPT-4 (2023). \texttt{arXiv:2303.12712}
  \item AgentBench: Evaluating LLMs as Agents (2023). \texttt{arXiv:2308.03688}
  \item Llama 2: Open Foundation and Fine-Tuned Chat Models (2023). \texttt{arXiv:2307.09288}
  \item ToolLLM: Facilitating Large Language Models to Master 16000+ Real-world
  APIs (2023). \texttt{arXiv:2307.16789}
  \item HuggingGPT: Solving AI Tasks with ChatGPT and its Friends in Hugging
  Face (2023). \texttt{arXiv:2303.17580}
  \item Self-Refine: Iterative Refinement with Self-Feedback (2023). \texttt{arXiv:2303.17651}
  \item CAMEL: Communicative Agents for "Mind" Exploration of Large Language
  Model Society (2023). \texttt{arXiv:2303.17760}
  \item Gorilla: Large Language Model Connected with Massive APIs (2023). \texttt{arXiv:2305.15334}
  \item Voyager: An Open-Ended Embodied Agent with Large Language Models (2023). \texttt{arXiv:2305.16291}
  \item WebArena: A Realistic Web Environment for Building Autonomous Agents (2023). \texttt{arXiv:2307.13854}
  \item Mind2Web: Towards a Generalist Agent for the Web (2023). \texttt{arXiv:2306.06070}
  \item MetaGPT: Meta Programming for A Multi-Agent Collaborative Framework (2023). \texttt{arXiv:2308.00352}
  \item The Rise and Potential of Large Language Model Based Agents: A Survey (2023). \texttt{arXiv:2309.07864}
  \item Judging LLM-as-a-Judge with MT-Bench and Chatbot Arena (2023). \texttt{arXiv:2306.05685}
  \item Language Models can Solve Computer Tasks (2023). \texttt{arXiv:2303.17491}
  \item Improving Factuality and Reasoning in Language Models through Multiagent
  Debate (2023). \texttt{arXiv:2305.14325}
  \item ChatDev: Communicative Agents for Software Development (2023). \texttt{arXiv:2307.07924}
  \item Describe, Explain, Plan and Select: Interactive Planning with Large
  Language Models Enables Open-World Multi-Task Agents (2023). \texttt{arXiv:2302.01560}
  \item Tool Learning with Foundation Models (2023). \texttt{arXiv:2304.08354}
  \item API-Bank: A Comprehensive Benchmark for Tool-Augmented LLMs (2023). \texttt{arXiv:2304.08244}
  \item AutoGen: Enabling Next-Gen LLM Applications via Multi-Agent Conversation (2023). \texttt{arXiv:2308.08155}
  \item Teaching Large Language Models to Self-Debug (2023). \texttt{arXiv:2304.05128}
  \item GAIA: a benchmark for General AI Assistants (2023). \texttt{arXiv:2311.12983}
  \item Augmented Language Models: a Survey (2023). \texttt{arXiv:2302.07842}
  \item Encouraging Divergent Thinking in Large Language Models through
  Multi-Agent Debate (2023). \texttt{arXiv:2305.19118}
  \item Visual ChatGPT: Talking, Drawing and Editing with Visual Foundation
  Models (2023). \texttt{arXiv:2303.04671}
  \item A Survey of Large Language Models (2023). \texttt{arXiv:2303.18223}
  \item Reasoning with Language Model is Planning with World Model (2023). \texttt{arXiv:2305.14992}
  \item ToolkenGPT: Augmenting Frozen Language Models with Massive Tools via
  Tool Embeddings (2023). \texttt{arXiv:2305.11554}
  \item ToolAlpaca: Generalized Tool Learning for Language Models with 3000
  Simulated Cases (2023). \texttt{arXiv:2306.05301}
  \item TaskMatrix.AI: Completing Tasks by Connecting Foundation Models with
  Millions of APIs (2023). \texttt{arXiv:2303.16434}
  \item AgentVerse: Facilitating Multi-Agent Collaboration and Exploring
  Emergent Behaviors (2023). \texttt{arXiv:2308.10848}
  \item MLAgentBench: Evaluating Language Agents on Machine Learning
  Experimentation (2023). \texttt{arXiv:2310.03302}
  \item Code Llama: Open Foundation Models for Code (2023). \texttt{arXiv:2308.12950}
  \item PaLM-E: An Embodied Multimodal Language Model (2023). \texttt{arXiv:2303.03378}
  \item LLM+P: Empowering Large Language Models with Optimal Planning
  Proficiency (2023). \texttt{arXiv:2304.11477}
  \item Plan-and-Solve Prompting: Improving Zero-Shot Chain-of-Thought Reasoning
  by Large Language Models (2023). \texttt{arXiv:2305.04091}
  \item ChatEval: Towards Better LLM-based Evaluators through Multi-Agent Debate (2023). \texttt{arXiv:2308.07201}
  \item Chameleon: Plug-and-Play Compositional Reasoning with Large Language
  Models (2023). \texttt{arXiv:2304.09842}
  \item On the Tool Manipulation Capability of Open-source Large Language Models (2023). \texttt{arXiv:2305.16504}
  \item Gemini: A Family of Highly Capable Multimodal Models (2023). \texttt{arXiv:2312.11805}
  \item Ghost in the Minecraft: Generally Capable Agents for Open-World
  Environments via Large Language Models with Text-based Knowledge and Memory (2023). \texttt{arXiv:2305.17144}
  \item Graph of Thoughts: Solving Elaborate Problems with Large Language Models (2023). \texttt{arXiv:2308.09687}
  \item ART: Automatic multi-step reasoning and tool-use for large language
  models (2023). \texttt{arXiv:2303.09014}
  \item RestGPT: Connecting Large Language Models with Real-World RESTful APIs (2023). \texttt{arXiv:2306.06624}
  \item Tool Learning in the Wild: Empowering Language Models as Automatic Tool
  Agents (2024). \texttt{arXiv:2405.16533}
  \item Large Language Model based Multi-Agents: A Survey of Progress and
  Challenges (2024). \texttt{arXiv:2402.01680}
  \item Exploring Large Language Model based Intelligent Agents: Definitions,
  Methods, and Prospects (2024). \texttt{arXiv:2401.03428}
  \item Large Language Model-Based Agents for Software Engineering: A Survey (2024). \texttt{arXiv:2409.02977}
  \item Small LLMs Are Weak Tool Learners: A Multi-LLM Agent (2024). \texttt{arXiv:2401.07324}
  \item Exploring Autonomous Agents through the Lens of Large Language Models: A
  Review (2024). \texttt{arXiv:2404.04442}
  \item LLM With Tools: A Survey (2024). \texttt{arXiv:2409.18807}
  \item Learning to Use Tools via Cooperative and Interactive Agents (2024). \texttt{arXiv:2403.03031}
  \item The Llama 3 Herd of Models (2024). \texttt{arXiv:2407.21783}
  \item \$\textbackslash{}tau\$-bench: A Benchmark for Tool-Agent-User Interaction in Real-World
  Domains (2024). \texttt{arXiv:2406.12045}
  \item MLE-bench: Evaluating Machine Learning Agents on Machine Learning
  Engineering (2024). \texttt{arXiv:2410.07095}
  \item Large Language Model Agent: A Survey on Methodology, Applications and
  Challenges (2025). \texttt{arXiv:2503.21460}
  \item From Language to Action: A Review of Large Language Models as Autonomous Agents and Tool Users (2025). \texttt{arXiv:2508.17281}
  \item AutoTool: Efficient Tool Selection for Large Language Model Agents (2025). \texttt{arXiv:2511.14650}
  \item Large Language Models for Agent-Based Modelling: Current and possible uses across the modelling cycle (2025). \texttt{arXiv:2507.05723}
  \item Improving Large Language Models Function Calling and Interpretability via Guided-Structured Templates (2025). \texttt{arXiv:2509.18076}
  \item Generalizability of Large Language Model-Based Agents: A Comprehensive Survey (2025). \texttt{arXiv:2509.16330}
  \item Fundamentals of Building Autonomous LLM Agents (2025). \texttt{arXiv:2510.09244}
  \item Position: Agent Should Invoke External Tools ONLY When Epistemically Necessary (2025). \texttt{arXiv:2506.00886}
  \item DeepSeek-R1: Incentivizing Reasoning Capability in LLMs via Reinforcement Learning (2025). \texttt{arXiv:2501.12948}
  \item Agentic Tool Use in Large Language Models (2026). \texttt{arXiv:2604.00835}
  \item Agentic Reasoning for Large Language Models (2026). \texttt{arXiv:2601.12538}
  \item The Evolution of Tool Use in LLM Agents: From Single-Tool Call to Multi-Tool Orchestration (2026). \texttt{arXiv:2603.22862}
  \item SAGE: Tool-Augmented LLM Task Solving Strategies in Scalable Multi-Agent Environments (2026). \texttt{arXiv:2601.09750}
  \item Benchmark Everything Everywhere All at Once (2026). \texttt{arXiv:2606.06462}
  \item What Twelve LLM Agent Benchmark Papers Disclose About Themselves: A Pilot Audit and an Open Scoring Schema (2026). \texttt{arXiv:2605.21404}
  \item ActuBench: A Multi-Agent LLM Pipeline for Generation and Evaluation of Actuarial Reasoning Tasks (2026). \texttt{arXiv:2604.20273}
  \item Agent Planning Benchmark: A Diagnostic Framework for Planning Capabilities in LLM Agents (2026). \texttt{arXiv:2606.04874}
\end{enumerate}

\end{document}
